\documentclass[../../../include/open-logic-section]{subfiles}

\begin{document}

\olfileid{sth}{z}{milestone}
\olsection{$\Z^-$：一座里程碑}

我们将在下一节再讨论 \stagesinf{}。然而，引入无穷公理后，我们便达到了一个重要的里程碑。现在我们已具备构成理论 $\Zminus$ 所需的全部公理。具体如下：

\begin{defn}
理论 $\Zminus$ 包含以下公理：外延公理、并集公理、对集公理、幂集公理、无穷公理，以及分离公理模式的所有实例。
\end{defn}

该名称意为 \emph{Zermelo}（策梅洛）集合论（\emph{减去}了某些我们稍后将提及的内容）。这一荣誉当属 Zermelo，因为他在其 \citeyear{Zermelo1908Untersuchungen} 中实质上提出了这一理论。\footnote{关于历史与技术的精彩评论，参见 \citet[附录 A]{Potter2004}。}

这一理论足够强大，足以让我们开展大量的数学工作。具体而言，你\emph{应当}回顾 \olref[sfr][][]{part}，并确信我们此前以非形式化方式所做的一切，都可以在~$\Zminus$ 内以更形式化的方式完成。（稍作尝试之后，你可能会想跳到前面读一下 \olref[sth][z][arbintersections]{sec}。）因此，从今往后，我们将默认（至少）在 $\Zminus$ 中工作，不再另行说明。

\end{document}